\chapter{Nền tảng và công nghệ sử dụng}

\section{Định hướng công nghệ của hệ thống}

Fitty được xây dựng như một ứng dụng Android hiện đại, trong đó giao diện, điều hướng, xử lý nghiệp vụ, lưu trữ cục bộ và dịch vụ phía sau được tách thành các lớp tương đối độc lập. Việc lựa chọn công nghệ không chỉ nhằm hiện thực đầy đủ chức năng mà còn phục vụ ba mục tiêu chính: dễ mở rộng, ổn định khi sử dụng thực tế và phù hợp với mô hình phát triển ứng dụng di động hiện nay \cite{android_jetpack}.

Về tổng thể, hệ thống sử dụng Kotlin làm ngôn ngữ phát triển chính, Jetpack Compose cho giao diện, Navigation Compose cho điều hướng, Hilt cho quản lý phụ thuộc, Room và DataStore cho lưu trữ cục bộ, Firebase cho các dịch vụ phía sau, WorkManager cho tác vụ nền, cùng các thư viện hỗ trợ media và truy cập dữ liệu ngoài khi cần thiết. Tập công nghệ này phù hợp với một ứng dụng sức khỏe cá nhân hóa, nơi dữ liệu người dùng, kế hoạch tập và nội dung bài tập đều cần được tổ chức và truy xuất nhất quán.

\section{Ngôn ngữ và công nghệ giao diện}

\subsection{Kotlin và Android}

Kotlin được chọn làm ngôn ngữ phát triển chính vì có cú pháp gọn, hỗ trợ tốt cho lập trình bất đồng bộ và tương thích cao với hệ sinh thái Jetpack \cite{android_jetpack}. Trên nền Android, lựa chọn này giúp hệ thống thuận lợi hơn trong việc tổ chức kiến trúc nhiều tầng, xử lý dữ liệu người dùng và mở rộng chức năng về sau.

\subsection{Jetpack Compose và Navigation Compose}

Jetpack Compose được dùng để xây dựng giao diện theo mô hình khai báo, phù hợp với đặc thù nhiều màn hình giàu trạng thái của Fitty \cite{android_compose}. Navigation Compose được sử dụng để tổ chức luồng điều hướng giữa các giai đoạn truy cập ban đầu, onboarding và vùng chức năng chính \cite{android_navigation_compose}. Hai công nghệ này giúp giao diện nhất quán hơn và làm rõ mối liên hệ giữa trạng thái màn hình với dữ liệu nghiệp vụ.

\subsection{Material 3}

Material 3 đóng vai trò là nền tảng thiết kế giao diện chung của ứng dụng. Việc sử dụng hệ thiết kế này giúp các thành phần hiển thị đồng nhất hơn, phù hợp với ứng dụng cần cân bằng giữa tính trực quan, dễ sử dụng và khả năng hiển thị dữ liệu sức khỏe rõ ràng.

\section{Kiến trúc và quản lý phụ thuộc}

\subsection{Kiến trúc nhiều tầng}

Hệ thống được tổ chức theo các tầng chính gồm giao diện, nghiệp vụ và dữ liệu. Cách phân lớp này giúp phần giao diện tập trung vào hiển thị và tương tác, trong khi quy tắc xử lý và truy cập dữ liệu được tách riêng để dễ bảo trì hơn. Đây là lựa chọn phù hợp với ứng dụng có nhiều phân hệ như truy cập, onboarding, kế hoạch tập, workout session và theo dõi tiến độ.

\subsection{Hilt}

Hilt được dùng để cấu hình dependency injection trong toàn bộ hệ thống \cite{android_hilt}. Công nghệ này giúp giảm sự phụ thuộc trực tiếp giữa các tầng, làm rõ vòng đời của các thành phần dùng chung và thuận tiện hơn khi kiểm thử hoặc thay thế từng phần của hệ thống.

\section{Lưu trữ và đồng bộ dữ liệu}

\subsection{Room và DataStore}

Room được sử dụng làm kho dữ liệu cục bộ cho các nội dung cần truy cập nhanh hoặc cần hỗ trợ ngoại tuyến, đặc biệt là metadata bài tập và một số dữ liệu phục vụ trải nghiệm hằng ngày \cite{android_room}. DataStore được dùng cho các cấu hình nhẹ và thông tin trạng thái phiên, phù hợp hơn với dữ liệu đơn giản trên thiết bị \cite{android_datastore}.

\subsection{Firebase}

Firebase cung cấp các dịch vụ phía sau chính của hệ thống, bao gồm xác thực người dùng, lưu trữ dữ liệu nghiệp vụ, lưu trữ media và thông báo đẩy \cite{firebase_setup,firebase_auth,firebase_firestore,firebase_storage,firebase_fcm}. Sự kết hợp này phù hợp với bài toán ứng dụng di động cần vừa quản lý dữ liệu người dùng, vừa hỗ trợ cập nhật nội dung động và duy trì tương tác ngoài phạm vi phiên làm việc hiện tại.

\subsection{Mô hình offline-first cho bài tập}

Đối với phân hệ bài tập, Fitty áp dụng mô hình offline-first: dữ liệu được lấy từ nguồn từ xa, chuẩn hóa, lưu xuống kho cục bộ rồi mới phục vụ giao diện. Cách làm này giúp giảm phụ thuộc vào trạng thái mạng tại thời điểm sử dụng, đồng thời cải thiện tốc độ mở danh mục và chi tiết bài tập. Việc tách metadata và media cũng giúp hệ thống tối ưu hơn về tài nguyên thiết bị và thời gian tải nội dung.

\section{Tác vụ nền, media và tích hợp ngoài}

WorkManager được sử dụng cho các tác vụ nền như đồng bộ dữ liệu định kỳ, nhờ đó giảm phụ thuộc vào thao tác thủ công của người dùng và hạn chế ảnh hưởng trực tiếp lên giao diện \cite{android_background_tasks,android_workmanager}. Về hiển thị nội dung trực quan, hệ thống sử dụng Coil cho ảnh và Media3/ExoPlayer cho video hướng dẫn bài tập \cite{android_media3}. Ngoài ra, Retrofit và các thư viện HTTP liên quan là cơ sở cho khả năng tích hợp dữ liệu ngoài trong tương lai \cite{retrofit_intro,retrofit_declarations}.

\section{Tổng kết chương}

Chương này đã trình bày ngắn gọn các công nghệ chính được sử dụng để xây dựng Fitty. Trọng tâm của việc lựa chọn công nghệ nằm ở khả năng tổ chức ứng dụng theo kiến trúc rõ ràng, hỗ trợ lưu trữ cục bộ và đồng bộ dữ liệu ổn định, đồng thời tạo nền tảng phù hợp cho việc mở rộng chức năng trong tương lai. Trên cơ sở đó, chương tiếp theo trình bày phần cài đặt, kiểm thử và đánh giá hệ thống.
